\documentclass[11pt,a4paper]{article}
\usepackage[italian]{babel}
\usepackage[margin=2.5cm]{geometry}
\usepackage{parskip}
\usepackage{amsmath, amssymb, amsthm}
\usepackage{booktabs}
\usepackage{array}
\usepackage{xcolor}
\usepackage{needspace}
\usepackage{enumitem}
\usepackage{listings}
\usepackage{tikz}
\usepackage[most]{tcolorbox}
\usepackage[hidelinks]{hyperref}
\usetikzlibrary{decorations.pathreplacing, shapes.geometric, arrows.meta, positioning}

% Niente vedove/orfane: i paragrafi non lasciano righe isolate a fine/inizio pagina
\widowpenalty=10000
\clubpenalty=10000

% Elenchi più compatti
\setlist{itemsep=2pt, parsep=0pt, topsep=4pt, partopsep=0pt}

% --- Riquadri con bordo nero, senza numero (si spezzano tra due pagine) ---
% Uso: \begin{definizione} ... \end{definizione}
%      \begin{teorema}[Nome] ... \end{teorema}  (il nome è facoltativo)
\tcbset{nota/.style={enhanced jigsaw, breakable, colback=white,
  colframe=black, boxrule=0.5pt, sharp corners}}
\NewTColorBox{definizione}{}{nota,
  before upper={\textbf{Definizione.}\ }}
\NewTColorBox{teorema}{o}{nota,
  before upper={\textbf{Teorema\IfValueT{#1}{ (#1)}.}\ }}
\NewTColorBox{proposizione}{o}{nota,
  before upper={\textbf{Proposizione\IfValueT{#1}{ (#1)}.}\ }}
\NewTColorBox{proprieta}{o}{nota,
  before upper={\textbf{Proprietà\IfValueT{#1}{ (#1)}.}\ }}
\NewTColorBox{assioma}{o}{nota, boxrule=1pt,
  before upper={\textbf{Assioma\IfValueT{#1}{ (#1)}.}\ }}

% --- Ambienti semplici (senza riquadro) ---
% Il titolo sta in cima al blocco, anche se il contenuto inizia con un riquadro o
% con codice; e non resta mai da solo in fondo a una pagina.
% Uso: \begin{esempio}[Titolo facoltativo] ... \end{esempio}
\newcommand{\newrunin}[2]{%
  \NewDocumentEnvironment{#1}{o}{%
    \par\Needspace{4\baselineskip}\addvspace{\medskipamount}\noindent
    \textbf{#2}\IfValueT{##1}{ (##1)}\textbf{.}\ \ignorespaces}%
  {\par\addvspace{\medskipamount}}}
\newrunin{esempio}{Esempio}
\newrunin{esercizio}{Esercizio}
\newrunin{osservazione}{Osservazione}

% V e F nelle tavole di verità
\newcommand{\V}{\textsf{V}}
\newcommand{\F}{\textsf{F}}

% --- Codice C ---
% Uso: \begin{codice} ... \end{codice}      codice C numerato
%      \begin{comando} ... \end{comando}    comandi da terminale
%      \begin{risultato} ... \end{risultato} output di un programma
%      \lstinline|printf("ciao");|           codice dentro al testo
\lstdefinestyle{c}{
  language=C,
  basicstyle=\ttfamily\small,
  keywordstyle=\color{blue}\bfseries,
  commentstyle=\color{gray}\itshape,
  stringstyle=\color{red!70!black},
  numbers=left,
  numberstyle=\tiny\color{gray},
  numbersep=8pt,
  columns=flexible,
  keepspaces=true,
  breaklines=true,
  showstringspaces=false,
  tabsize=4,
  morekeywords={include, define},
  % lettere accentate nei commenti
  literate={à}{{\`a}}1 {è}{{\`e}}1 {é}{{\'e}}1 {ì}{{\`i}}1
           {ò}{{\`o}}1 {ù}{{\`u}}1 {È}{{\`E}}1 {À}{{\`A}}1
}
\lstset{style=c}
\newtcblisting{codice}{listing only, nota, left=6mm,
  listing options={style=c}}
\newtcblisting{comando}{listing only, nota,
  listing options={basicstyle=\ttfamily\small, numbers=none, language={},
    columns=flexible, keepspaces=true, breaklines=true}}
\newtcblisting{risultato}{listing only, enhanced jigsaw, breakable,
  colback=black!5, colframe=black, boxrule=0.5pt, sharp corners,
  listing options={basicstyle=\ttfamily\small, numbers=none, language={},
    columns=flexible, keepspaces=true, breaklines=true}}

% --- Diagrammi di flusso (simboli standard) ---
\tikzset{
  terminale/.style = {ellipse, draw, minimum width=2.5cm, minimum height=0.9cm},
  io/.style        = {trapezium, trapezium left angle=70, trapezium right angle=110,
                      draw, minimum width=2.5cm, minimum height=0.9cm},
  processo/.style  = {rectangle, draw, minimum width=2.5cm, minimum height=0.9cm},
  decisione/.style = {diamond, draw, aspect=2, minimum width=2.5cm},
  freccia/.style   = {thick, -{Stealth}}
}

\title{Invertibili e divisori dello zero modulo $n$}
\author{Marco Cerbella}
\date{Lezione 13 -- 2 novembre 2026}

\begin{document}
\maketitle

\section{Divisibilità, invertibili e divisori dello zero in $\mathbb{Z}_n$}

\begin{definizione}
Siano $[a]_n, [b]_n \in \mathbb{Z}_n$.
\begin{itemize}
  \item Diciamo che $[a]_n$ \emph{divide} $[b]_n$ ($[a]_n \mid [b]_n$) se $\exists [c]_n \in \mathbb{Z}_n$ tale che $[b]_n = [a]_n\cdot[c]_n$.
  \item $[a]_n \in \mathbb{Z}_n$ è \emph{invertibile} se $\exists [a']_n \in \mathbb{Z}_n$ tale che $[a]_n\cdot[a']_n = [1]_n$.
  \item $[a]_n \in \mathbb{Z}_n$, $[a]_n \neq [0]_n$, è un \emph{divisore dello zero} se $\exists [b]_n \neq [0]_n$ tale che $[a]_n\cdot[b]_n = [0]_n$.
\end{itemize}
\end{definizione}

\begin{osservazione}
In generale $\mathbb{Z}_n$ non è un dominio di integrità, cioè in $\mathbb{Z}_n$ non vale la legge di annullamento del prodotto: ad esempio $[2]_6\cdot[3]_6 = [6]_6 = [0]_6$.

Non sempre vale la legge di cancellazione del prodotto. Ad esempio in $\mathbb{Z}_{12}$: $[5]\cdot[8] = [40] = [4]$ e $[8]\cdot[8] = [64] = [4]$ ($64 = 12\cdot 5 + 4$), quindi $[5]\cdot[8] = [8]\cdot[8]$ ma $[5] \neq [8]$.

Però $\mathbb{Z}_2$, $\mathbb{Z}_3$, $\mathbb{Z}_5$ sono campi.
\end{osservazione}

\section{Proprietà delle congruenze}

\begin{proprieta}[Corollario]
Sia $n > 0$ e $\forall a, b, c \in \mathbb{Z}$. Se $a \equiv b \bmod n$, allora
\[
  a + c \equiv b + c \bmod n, \qquad a\cdot c \equiv b\cdot c \bmod n, \qquad a^i \equiv b^i \bmod n .
\]
In termini di classi: $[a]_n = [b]_n \Rightarrow [a + c]_n = [b + c]_n$, $[ac]_n = [bc]_n$, $[a^i]_n = [b^i]_n$.
\end{proprieta}

\begin{proprieta}
\begin{itemize}
  \item $ac \equiv bc \bmod n$ e $\operatorname{MCD}(c, n) = 1 \Rightarrow a \equiv b \bmod n$ (cioè $[ac]_n = [bc]_n$ e $\operatorname{MCD}(c, n) = 1 \Rightarrow [a]_n = [b]_n$).
  \item $ac \equiv bc \bmod n$ e $\operatorname{MCD}(c, n) = d \Rightarrow a\cdot\frac{c}{d} \equiv b\cdot\frac{c}{d} \bmod \frac{n}{d} \Rightarrow a \equiv b \bmod \frac{n}{d}$.
  \item $a \equiv b \bmod n$ e $p \mid n \Rightarrow a \equiv b \bmod p$.
  \item $a \equiv b \bmod r$ e $a \equiv b \bmod s \Rightarrow a \equiv b \bmod \operatorname{mcm}(r, s)$.
\end{itemize}
\end{proprieta}

\begin{osservazione}[Rappresentanti di $\mathbb{Z}_n$]
Si può rappresentare $\mathbb{Z}_n$ con qualunque insieme di $n$ interi consecutivi. Ad esempio
\[
  \mathbb{Z}_6 = \{[0], [1], [2], [3], [4], [5]\} = \{[7], [8], [9], [10], [11], [12]\} = \{[-2], [-1], [0], [1], [2], [3]\},
\]
in generale $\mathbb{Z}_n = \{[c]_n, [c + 1]_n, \dots, [c + (n-1)]_n\}$. Per esempio in $\mathbb{Z}_6$: $[0]_6 = \{6k \mid k \in \mathbb{Z}\} = [6]_6 = [12]_6 = \{\dots, -6, 0, 6, 12, \dots\}$; $[1]_6 = \{6k + 1\} = [7]_6 = \{\dots, -11, -5, 1, 7, 13, \dots\}$; $[2]_6 = \{6k + 2\} = [8]_6 = [-4]_6 = \{\dots, -10, -4, 2, 8, \dots\}$.
\end{osservazione}

\begin{proposizione}
Sia $c$ tale che $\operatorname{MCD}(c, n) = 1$. Allora
\[
  \mathbb{Z}_n = \{[c\cdot 0], [c\cdot 1], [c\cdot 2], \dots, [c(n-1)]\} .
\]
\end{proposizione}

\begin{esempio}
In $\mathbb{Z}_6$, con $c = 5$ e $\operatorname{MCD}(5, 6) = 1$:
\[
  \mathbb{Z}_6 = \{[5\cdot 0], [5\cdot 1], [5\cdot 2], [5\cdot 3], [5\cdot 4], [5\cdot 5]\} = \{[0], [5], [10], [15], [20], [25]\}.
\]
\end{esempio}

\section{Invertibili e divisori dello zero}

\begin{proposizione}
Sia $[a]_n \in \mathbb{Z}_n$, $[a]_n \neq [0]_n$. Allora
\[
  [a]_n \text{ è un divisore dello zero} \iff [a]_n \text{ non è invertibile.}
\]
\end{proposizione}

\begin{proof}
$\Rightarrow$ Sia $[a]$ un divisore dello zero, cioè $\exists [b] \neq [0]$ tale che $[a]\cdot[b] = [0]$. Se $[a]$ fosse invertibile, esisterebbe $[a']$ con $[a]\cdot[a'] = [1]$ e
\[
  [a']\cdot[a]\cdot[b] = [a']\cdot[0] \ \Rightarrow\ [1]\cdot[b] = [b] = [0] ,
\]
ma $[b] \neq [0]$: assurdo.

$\Leftarrow$ Sia $[a]$ non invertibile, cioè $\forall [b] \in \{[0], \dots, [n-1]\}$ si ha $[a]\cdot[b] \neq [1]$. Allora gli $n$ prodotti $[a]\cdot[b]$ assumono al più $n - 1$ valori distinti, quindi devono esserci due classi $[b_1] \neq [b_2]$ tali che $[a]\cdot[b_1] = [a]\cdot[b_2]$, cioè $[ab_1] = [ab_2]$, cioè $ab_1 \equiv ab_2 \bmod n$, cioè $\exists k \in \mathbb{Z}$ tale che $ab_1 - ab_2 = kn$. Cioè $a(b_1 - b_2) \equiv 0 \bmod n$, cioè $[a]\cdot[b_1 - b_2] = [0]$, con $[a] \neq [0]$ per ipotesi e $[b_1 - b_2] \neq [0]$ perché $b_1 \neq b_2$ (come classi). Quindi $[a]$ è un divisore dello zero.
\end{proof}

\begin{proposizione}
In $\mathbb{Z}_n$ la classe $[c]_n$ è invertibile se e solo se $\operatorname{MCD}(c, n) = 1$.
\end{proposizione}

\begin{proof}
$\Rightarrow$ Sia $[c]$ invertibile: $\exists [c'] \in \mathbb{Z}_n$ con $[c]\cdot[c'] = [1]$, cioè $[c\cdot c'] = [1]$, cioè $c\cdot c' - 1 = kn$ per qualche $k$, cioè $c\cdot c' - kn = 1$. Allora $\operatorname{MCD}(c, n) \mid (c\cdot c' - kn) = 1$, cioè $\operatorname{MCD}(c, n) = 1$.

$\Leftarrow$ Sia $\operatorname{MCD}(c, n) = 1$. Per il teorema di Bézout esistono $\alpha, \beta \in \mathbb{Z}$ tali che $1 = \operatorname{MCD}(c, n) = \alpha c + \beta n$, cioè $\alpha c - 1 = -\beta n$, cioè $\alpha c \equiv 1 \bmod n$, cioè $[\alpha]_n\cdot[c]_n = [1]_n$. Quindi $[\alpha]$ è l'inverso di $[c]$ modulo $n$.
\end{proof}

\begin{esempio}[Inverso di $[13]_{30}$ in $\mathbb{Z}_{30}$]
Si usa l'algoritmo di Euclide \emph{esteso} tra $13$ e $30$. Si costruisce la tabella con colonne $q$, $r$, $s$, $t$ (si parte con $r = 30, 13$; $s = 1, 0$; $t = 0, 1$; e ad ogni passo $q$ è il quoziente, $r$ il resto, $s_{\text{nuovo}} = s_{\text{prec-prec}} - q\cdot s_{\text{prec}}$, e analogamente per $t$):
\[
\begin{array}{c|rrr}
  q & r & s & t \\ \hline
  - & 30 & 1 & 0 \\
  - & 13 & 0 & 1 \\
  2 & 4 & 1 & -2 \\
  3 & 1 & -3 & 7 \\
  4 & 0 & 13 & -30
\end{array}
\]
L'ultimo resto non nullo è $1 = \operatorname{MCD}(30, 13)$, con $s = -3$ e $t = 7$:
\[
  1 = 30\cdot(-3) + 13\cdot 7 .
\]
Quindi $[7]_{30}$ è l'inverso di $[13]_{30}$. Verifica: $[7]\cdot[13] = [91] = [1]_{30}$, perché $91 = 3\cdot 30 + 1$.
\end{esempio}

\end{document}
